% ============================================================
% CAPITOLO 24 - RITARDI COMPLESSI
% ============================================================
\chapter{Ritardi complessi}
\label{chap:ritardi_complessi}

Oltre al 4-3 si usano altri quattro ritardi, che coinvolgono la nona e la settima. Nella cifra \textbf{X--Y} il primo numero indica l'intervallo del ritardo sul basso e il secondo quello della risoluzione.

\begin{center}
\begin{tabularx}{\linewidth}{lLLL}
\toprule
\textbf{Ritardo} & \textbf{Preparazione} & \textbf{Nota ritardata} & \textbf{Risoluzione} \\
\midrule
9-8 & V (quinta del V) & nona sul basso & fondamentale: tonica in stato fondamentale \\
9-3 & IV, iv (terza) & nona sul V & terza di una tonica in primo rivolto (il basso scende di terza) \\
7-6 & V (quinta del V) & settima sul basso & sesta: tonica in primo rivolto \\
7-6 sul II grado & I (fondamentale) & settima sul II grado & sesta \emph{maggiore}, la sensibile: \sigla{VII}{6}{} (triade diminuita in primo rivolto); in minore si cifra $\sharp$6 \\
\bottomrule
\end{tabularx}
\end{center}

\begin{regola}[Condotta dei ritardi complessi]
	\label{reg:condotta_ritardo}
	Valgono le regole del ritardo semplice (Regola \ref{reg:condotta_ritardo_base}), con queste precisazioni:
	\begin{itemize}
		\item la nota di risoluzione non compare in nessuna voce superiore durante il ritardo; fa eccezione la risoluzione sulla stessa nota del basso, come nel 9-8 (il Do è già nel basso);
		\item la \textbf{sensibile} del V sale alla tonica. Quando la tonica è la nota di risoluzione del ritardo (9-8, 7-6) la sensibile non può salire durante il ritardo: se è in una \textbf{voce interna} può scendere di terza sulla quinta della tonica (eccezione di Bach, Osservazione \ref{oss:eccezione_bach}); nel soprano non può mai scendere.
	\end{itemize}
\end{regola}

\section{I ritardi 9-8 e 9-3}
\label{sec:ritardi_nona}

\begin{esempio}[Ritardo 9-8 nel soprano]
	\label{ex:ritardo_9_8_s}
	La nona Re (quinta del V) è tenuta sul Do del basso e scende al Do. La sensibile, in una voce interna, cade di terza sulla quinta della tonica: sul Sol in Do maggiore (nel tenore), sul Mi in La minore (nel contralto).
	\musica{es_cap10_9_8_s}
\end{esempio}

\begin{esempio}[Ritardo 9-8 nel contralto]
	\label{ex:ritardo_9_8_a}
	Il ritardo del contralto in 9-8 è realizzabile solo in Do maggiore e in posizione stretta. In La minore non esiste una realizzazione corretta a quattro voci: la sensibile non può salire (comparirebbe il La contro il Si ritardato), nel soprano non può cadere, e nel tenore la caduta di terza obbligherebbe a una distanza eccessiva tra le voci.
	\musica{es_cap10_9_8_a}
\end{esempio}

\begin{esempio}[Ritardo 9-3]
	\label{ex:ritardo_9_3}
	Nel passaggio dal V alla tonica in primo rivolto il basso scende di terza (Sol-Mi, Mi-Do). In La minore la nona sopra il Mi è il Fa, una nona minore; la cifra resta 9, perché il Fa appartiene già alla scala e non va alterato. Ritardo nel soprano:
	\musica{es_cap10_9_3_s}
	Ritardo nel contralto:
	\musica{es_cap10_9_3_a}
\end{esempio}

\section{I ritardi 7-6 sulla tonica e sul II grado}
\label{sec:ritardi_settima}

\begin{esempio}[Ritardo 7-6]
	\label{ex:ritardo_7_6}
	La settima (Re sul Mi del basso) scende alla sesta (Do): la risoluzione è un accordo di sesta. Nel soprano:
	\musica{es_cap10_7_6_s}
	Il ritardo 7-6 nel contralto non si realizza a quattro voci nel rispetto delle spaziature: la sensibile può cadere solo nel tenore, e il soprano, per non produrre ottave parallele con il basso, dovrebbe staccarsi dal contralto di più di un'ottava.
\end{esempio}

\begin{esempio}[Ritardo 7-6 sul II grado]
	\label{ex:ritardo_7_6d}
	La settima Do è tenuta sul Re del basso e scende di semitono alla sensibile Si (sesta maggiore): l'accordo di risoluzione è \sigla{VII}{6}{}. In La minore la sesta è la sensibile alzata Sol$\sharp$ e la cifra diventa $\sharp$6. La sensibile sale alla tonica, la quinta diminuita Fa scende a Mi. Ritardo nel soprano:
	\musica{es_cap10_7_6d_s}
	Ritardo nel contralto:
	\musica{es_cap10_7_6d_a}
\end{esempio}
